\documentclass[11pt]{article}
\usepackage[a4paper,margin=2.3cm]{geometry}
\usepackage[utf8]{inputenc}
\usepackage[T1]{fontenc}
\usepackage[brazilian]{babel}
\usepackage{hyperref}
\usepackage{longtable}
\usepackage{booktabs}
\usepackage{array}
\usepackage{enumitem}
\usepackage{xcolor}

\definecolor{codebg}{RGB}{245,245,245}
\usepackage{listings}
\lstset{
  backgroundcolor=\color{codebg},
  basicstyle=\ttfamily\small,
  breaklines=true,
  frame=single,
  columns=fullflexible,
  showstringspaces=false,
}

\hypersetup{colorlinks=true, linkcolor=blue, urlcolor=blue, citecolor=blue}
\setlength{\parskip}{0.6em}
\setlength{\parindent}{0pt}

\newcolumntype{P}[1]{>{\raggedright\arraybackslash}p{#1}}

\title{\texttt{pipeline\_lab} --- Referência Completa\\
\large Funções, parâmetros e fundamentação na literatura}
\author{modelagem-lab}
\date{\today}

\begin{document}
\maketitle

\begin{abstract}
\noindent Este documento cataloga toda função pública de \texttt{pipeline\_lab}
(\url{python/pipeline_lab/} no repositório), seus parâmetros, e a passagem de
literatura que justifica cada decisão de design relevante. A biblioteca é uma
coleção de funções soltas (não uma API fluente) para montar o funil de
modelagem de risco de crédito --- divisão, construção de variáveis, agregação
temporal, descoberta de interação, categorização/WOE, pré-seleção e a busca
stepwise Pedro\_Wise --- em cima de um \texttt{pandas.DataFrame} qualquer, sem
depender de disco ou de um servidor web. As citações vêm da base de
literatura já curada em \texttt{docs/literatura/} do repositório.
\end{abstract}

\tableofcontents
\newpage

\section{Filosofia e convenções}

\texttt{pipeline\_lab} é uma coleção de \textbf{funções soltas}, não uma API
orientada a objeto: cada etapa recebe DataFrames e devolve DataFrames novos,
nunca muta o original, nunca lê ou escreve em disco. Isso permite plugar um
\texttt{DataFrame} qualquer e escolher só os módulos que fazem sentido para o
caso de uso, na ordem:

\begin{lstlisting}
divisao -> construcao (opcional) -> esfera1 (opcional)
        -> esfera2 (opcional) -> categorizar
        -> preselecao (opcional) -> treinamento
\end{lstlisting}

\textbf{Convenção de coluna-alvo.} A partir de \texttt{divisao}, a variável
resposta deve se chamar \texttt{"y"} em \texttt{df\_dev}/\texttt{df\_teste} ---
todas as etapas seguintes acessam \texttt{df["y"]} diretamente.
\texttt{dividir\_por\_amostra}/\texttt{dividir\_aleatorio} aceitam
\texttt{coluna\_y} para fazer esse rename automaticamente.

\textbf{Anti-leakage por construção.} Toda etapa que "ajusta" algo a partir de
\texttt{y} (WOE, tabela de bins monotônicos, regras de interação) faz isso só
em dev, e reaplica a mesma tabela/regra em teste sem reajustar --- nunca
recalcula usando o \texttt{y} de teste. É o mesmo padrão \emph{fit/transform}
do scikit-learn.

\textbf{Convenção de "base".} Uma variável crua pode ter várias versões
(\texttt{idade\_woe}, \texttt{idade\_log}, \texttt{idade\_bin}) --- todas
compartilham o mesmo prefixo semântico ("base"). O Pedro\_Wise nunca deixa
duas versões da mesma base coexistirem no modelo final.

\section{\texttt{divisao}}

Divide um DataFrame em \texttt{(df\_dev, df\_teste)}. Sempre o primeiro passo.

\subsection*{\texttt{dividir\_por\_amostra(df, coluna\_amostra, valores\_dev, valores\_teste, coluna\_y=None)}}

Split por uma coluna de amostra que já existe no dataframe --- não assume
nomes fixos: você diz quais rótulos contam como dev e quais contam como
teste.

\begin{longtable}{P{3.2cm}P{2cm}P{2cm}P{7.5cm}}
\toprule
\textbf{Parâmetro} & \textbf{Tipo} & \textbf{Default} & \textbf{Explicação} \\
\midrule
\endhead
\texttt{df} & DataFrame & --- & Dataframe completo, ainda não dividido. \\
\texttt{coluna\_amostra} & str & --- & Coluna que identifica a amostra de cada linha (ex.\ \texttt{"split"}). Valores livres --- \texttt{"DES"/"OOT"}, \texttt{"train"/"val"}, o que já vier na fonte. \\
\texttt{valores\_dev} & list[str] & --- & Quais valores contam como desenvolvimento. \\
\texttt{valores\_teste} & list[str] & --- & Quais valores contam como teste. Linhas fora das duas listas são descartadas. \\
\texttt{coluna\_y} & str\,|\,None & None & Se passado, renomeia essa coluna para \texttt{"y"} antes de dividir. \\
\bottomrule
\end{longtable}

\subsection*{\texttt{dividir\_aleatorio(df, proporcao\_teste=0.5, semente=42, coluna\_y=None)}}

Split aleatório simples --- embaralha e corta.

\begin{longtable}{P{3.2cm}P{2cm}P{2cm}P{7.5cm}}
\toprule
\textbf{Parâmetro} & \textbf{Tipo} & \textbf{Default} & \textbf{Explicação} \\
\midrule
\endhead
\texttt{proporcao\_teste} & float & 0.5 & Fração das linhas que vai para teste. \\
\texttt{semente} & int & 42 & \texttt{random\_state} --- reprodutibilidade. \\
\texttt{coluna\_y} & str\,|\,None & None & Ver acima. \\
\bottomrule
\end{longtable}

\section{\texttt{construcao}}

Razões e diferenças interpretáveis entre pares de variáveis relacionadas ---
ex.\ \texttt{pago / fatura} = "proporção paga da fatura", uma feature de
negócio óbvia que não existe nas colunas originais.

Escopo deliberadamente mínimo: não é busca automática tipo Deep Feature
Synthesis ou programação genética. A escolha de v1 foi gerar candidatas
simples e interpretáveis primeiro --- mais barato e testável imediatamente
contra o Pedro\_Wise --- em vez da escola "automática/exaustiva" descrita no
paper fundador de Kanter \& Veeramachaneni (2015), \emph{Deep
Feature Synthesis: Towards automating data science endeavors}, que aplica
operações de agregação recursivamente sobre relações entre tabelas (base da
biblioteca \texttt{featuretools}).

Um paralelo direto em crédito: Dumitrescu et al.\ (2021), \emph{Machine
learning for credit scoring: Improving logistic regression with non-linear
decision-tree effects}, usa efeitos extraídos de árvores como features
construídas de volta para uma regressão logística --- o mesmo espírito de
"construção assistida" que \texttt{esfera2} (RuleFit-style) aplica depois de
\texttt{construcao} no funil.

\begin{longtable}{P{4cm}P{9.5cm}}
\toprule
\textbf{Função} & \textbf{Assinatura e parâmetros} \\
\midrule
\endhead
\texttt{construir\_razao} &
\texttt{(numerador, denominador, nome, epsilon=1e-6)} --- \texttt{epsilon}
substitui denominador exatamente zero (mantém o sinal no resto, só evita
divisão por zero literal); comum em variáveis monetárias reais. \\[0.4em]
\texttt{construir\_diferenca} &
\texttt{(a, b, nome)} --- \texttt{a - b}, útil para variáveis na mesma
unidade (ex.\ variação de fatura mês a mês). \\[0.4em]
\texttt{construir\_razoes\_em\_lote} &
\texttt{(df, pares, epsilon=1e-6)} --- \texttt{pares} é uma lista de
\texttt{(numerador, denominador, nome)}, aplica \texttt{construir\_razao}
para cada tripla. Retorna só as colunas construídas. \\
\bottomrule
\end{longtable}

\section{\texttt{esfera1} --- agregação temporal}

Agregação "comportamental" (behavioral scoring): a partir de um painel
(chave + tempo + variável observada mês a mês), constrói primitivas de
janela móvel (máximo, média, mínimo, desvio-padrão, tendência) e devolve
uma linha por chave.

Catálogo de primitivas inspirado no vocabulário de agregação do Deep
Feature Synthesis (Kanter \& Veeramachaneni, 2015) --- sem
adotar a biblioteca inteira, só o vocabulário relevante para painéis de
crédito mensal.

\textbf{Garantia sem look-ahead}: cada linha de saída usa só observações até
e incluindo o próprio período --- nunca dados futuros da mesma chave.

\textbf{Preservação de split}: \texttt{aplicar} agrega \texttt{df\_dev} e
\texttt{df\_teste} \emph{separadamente} --- nunca junta as duas bases para
agregar e depois redividir, o que contaminaria a fronteira temporal do
split.

\begin{longtable}{P{4cm}P{9.5cm}}
\toprule
\textbf{Função} & \textbf{Assinatura e parâmetros} \\
\midrule
\endhead
\texttt{agregar} &
\texttt{(df, chave, coluna\_tempo, colunas\_valor, janelas)} --- core de
agregação sem split, uma chamada por base. \texttt{chave} identifica a
entidade; \texttt{coluna\_tempo} deve ser ordenável por chave;
\texttt{colunas\_valor} são as colunas numéricas agregadas;
\texttt{janelas} é uma lista de tamanhos de janela móvel em períodos
(ex.\ \texttt{[3, 6, 12]}). Retorna \texttt{(df\_agregado, colunas\_geradas)}. \\[0.4em]
\texttt{aplicar} &
\texttt{(df\_dev, df\_teste, chave, coluna\_tempo, colunas\_valor, janelas)}
--- chama \texttt{agregar} separadamente em dev e teste. Retorna
\texttt{(df\_dev\_agregado, df\_teste\_agregado, colunas\_geradas)}. \\
\bottomrule
\end{longtable}

\section{\texttt{esfera2} --- descoberta de interação}

Descoberta de interação RuleFit-style (Friedman \& Popescu, 2008,
\emph{Predictive learning via rule ensembles},
\emph{Annals of Applied Statistics}): treina um ensemble de árvores rasas
(\texttt{GradientBoostingClassifier}) sobre as candidatas já construídas,
extrai os caminhos raiz-folha como regras de interação (ex.\ "tendência
alta E severidade alta"), e devolve as regras como candidatas avaliáveis
para o funil.

\textbf{Diferença deliberada do RuleFit original}: no paper, as regras
extraídas viram termos de uma regressão com penalização L1 (Lasso) que
decide os pesos finais. Aqui as regras viram colunas 0/1 que entram no funil
normal --- quem decide o que fica no modelo final continua sendo a busca
stepwise do Pedro\_Wise, não uma regressão L1 própria.

\subsection*{\texttt{transformar\_categoricas\_woe(dev, teste, colunas\_x, colunas\_categoricas=None)}}

O \texttt{GradientBoostingClassifier} só aceita entrada numérica, então
colunas categóricas são WOE-codificadas antes de entrar na árvore. Colunas
de texto em \texttt{colunas\_x} são tratadas como categóricas
automaticamente. É um encoding "ingênuo" (sem regularização/shrinkage) ---
Pargent, Pfisterer, Thomas \& Bischl (2022), \emph{Regularized target
encoding outperforms traditional methods...}, mostra que versões
regularizadas de target encoding superam a versão ingênua, que pode vazar
informação do target e overfittar; ponto de atenção registrado, ainda não
implementado aqui (o WOE final que o modelo vê é ajustado só em dev, o
mesmo cuidado anti-leakage do resto do lab).

\subsection*{\texttt{aplicar(df\_dev, df\_teste, ...)}}

\begin{longtable}{P{4cm}P{2cm}P{1.6cm}P{6.5cm}}
\toprule
\textbf{Parâmetro} & \textbf{Tipo} & \textbf{Default} & \textbf{Explicação} \\
\midrule
\endhead
\texttt{colunas\_categoricas} & list[str]\,|\,None & None & Ver \texttt{transformar\_categoricas\_woe}. \\
\texttt{profundidade\_maxima} & int & 2 & Profundidade das árvores --- controla o tamanho máximo de uma regra (nº de condições). \\
\texttt{n\_arvores} & int & 60 & Nº de árvores no GBM --- mais árvores, mais candidatas de regra antes da deduplicação/filtro. \\
\texttt{min\_suporte} & float & 0.02 & Suporte mínimo (fração de linhas que satisfaz a regra) --- descarta regras quase nunca verdadeiras. \\
\texttt{max\_suporte} & float & 0.5 & Suporte máximo --- descarta regras quase sempre verdadeiras (constante disfarçada). \\
\texttt{max\_regras} & int & 10 & Teto de regras devolvidas após os filtros. \\
\texttt{permitir\_cruzamento\_entre\_bases} & bool & True & Se False, restringe regras a combinarem só primitivas da mesma variável bruta --- útil quando não se quer misturar domínios numa condição só. \\
\texttt{proporcao\_variaveis\_por\_split} & float\,|\,None & None & Fração de variáveis candidatas por split (equivalente a \texttt{max\_features} do sklearn) --- diversifica as regras extraídas. \\
\texttt{iv\_minimo} & float & 0.02 & Só regras com IV de \emph{teste} $\geq$ esse valor viram coluna --- nunca o IV de dev, que já foi usado para escolher a regra e está inflado por construção. \\
\bottomrule
\end{longtable}

Retorna \texttt{(df\_dev, df\_teste, colunas\_geradas)}.

\section{\texttt{categorizar} --- binning + WOE}

Sempre a última etapa antes da pré-seleção/treinamento --- é aqui que as
versões alternativas de cada variável (\texttt{\_woe}, \texttt{\_log},
\texttt{\_bin}) nascem. Esfera 2 vem antes justamente porque essas
transformações ainda não existem naquele ponto do funil.

\subsection*{\texttt{categorizar\_e\_transformar(df\_dev, df\_teste, gerar\_transformacoes\_potencia=True, gerar\_bin\_ordinal=True, ao\_processar\_coluna=None)}}

Para cada coluna (exceto \texttt{y}):

\begin{itemize}[leftmargin=1.2cm]
\item \textbf{Numérica contínua} (\texttt{nunique > 2}): binning monotônico
primeiro (\texttt{bins\_monotonicos} --- merge guloso que força a taxa de
evento a ser monotônica entre bins adjacentes; versão pragmática, não ótima
via MIP, da ideia central do OptBinning ---
Navas-Palencia, 2020,
\emph{Optimal binning: mathematical programming formulation},
hoje o padrão prático de mercado para binning em crédito).
\item \textbf{Categórica/binária} (texto, ou numérica com só 2 valores ---
ex.\ uma flag 0/1 de regra da Esfera 2): sem binning, cada valor já é o
"bin". \emph{Bug real corrigido}: colunas binárias eram descartadas
silenciosamente aqui antes da correção de 2026-07 (\texttt{bins\_monotonicos}
numa coluna de 2 valores gera só 2 edges, insuficiente para discretizar).
\end{itemize}

Depois do binning: \texttt{ajustar\_woe} (Siddiqi, \emph{Credit Risk
Scorecards} --- convenção de sinal $\mathrm{WOE} = \ln(\%\text{não-evento} /
\%\text{evento})$, positivo quando o bin tem mais "bons pagadores" relativo
à distribuição total) ajustado só em dev, reaplicado em teste.

\begin{longtable}{P{4.2cm}P{2cm}P{1.6cm}P{6cm}}
\toprule
\textbf{Parâmetro} & \textbf{Tipo} & \textbf{Default} & \textbf{Explicação} \\
\midrule
\endhead
\texttt{gerar\_transformacoes\_potencia} & bool & True & Gera também log/raiz/quad/cubo/inversas como candidatas extras --- o Pedro\_Wise testa trocar a versão WOE por uma dessas via a semântica de "base". Só numérica contínua; descarta transformações com domínio inválido. Análogo em espírito a Box--Cox/Yeo--Johnson (Atkinson, Riani \& Corbellini, 2021). \\
\texttt{gerar\_bin\_ordinal} & bool & True & Gera também o índice do bin (faixa) como candidata extra --- outra "versão" da mesma base. \\
\texttt{ao\_processar\_coluna} & Callable\,|\,None & None & Callback \texttt{(coluna, iv)} chamado após cada coluna processada --- gancho de progresso em tempo real, sem a biblioteca precisar saber o que é uma fila/WebSocket. \\
\bottomrule
\end{longtable}

Retorna \texttt{(woe\_dev, woe\_teste, iv\_por\_variavel)}.

\section{\texttt{preselecao}}

Reduz o volume de candidatas antes do Pedro\_Wise. Motivação: depois que
Construção e as transformações de potência de \texttt{categorizar} geram
dezenas/centenas de colunas extras por variável-base, o volume pode
explodir --- mais lento e maior risco de selecionar ruído por acaso (o
problema de múltiplos testes, o mesmo pano de fundo da literatura de
\emph{stability selection}: Meinshausen \& B\"uhlmann, 2008/2010,
\emph{Stability Selection}, JRSS-B).

Três filtros em sequência, cada um opcional (\texttt{limiar=None} pula):

\begin{longtable}{P{4cm}P{9.5cm}}
\toprule
\textbf{Função} & \textbf{Parâmetros e explicação} \\
\midrule
\endhead
\texttt{filtrar\_variancia(df, colunas, limiar=1e-6)} &
Mantém colunas com variância $>$ \texttt{limiar} --- descarta
quase-constantes, comuns entre transformações de potência quando a
variável original tem pouca variação. \\[0.4em]
\texttt{filtrar\_iv(colunas, iv\_por\_base, limiar=0.02)} &
Mantém colunas cuja variável-base tem IV $\geq$ \texttt{limiar}. Bases sem
entrada em \texttt{iv\_por\_base} são tratadas como IV=0 (descartadas). \\[0.4em]
\texttt{filtrar\_correlacao(df, colunas, iv\_por\_base, limiar=0.9)} &
Entre pares com $|r| \geq$ \texttt{limiar}, mantém o de maior IV.
\textbf{Pares da mesma base são ignorados de propósito} --- filtrar aqui
colapsaria a família inteira num só sobrevivente antes do Pedro\_Wise
escolher a melhor versão. \\[0.4em]
\texttt{pre\_selecionar(df, iv\_por\_base, limiar\_variancia=1e-6, limiar\_iv=0.02, limiar\_correlacao=0.9)} &
Aplica os 3 filtros em sequência. Retorna \texttt{colunas\_mantidas},
contagens em cada etapa do funil, e pares descartados por correlação. \\
\bottomrule
\end{longtable}

Este design --- filtrar correlação entre bases diferentes, nunca dentro da
mesma base --- evita o modo de falha descrito por Faletto \& Bien (2022),
\emph{Cluster Stability Selection}: quando há
proxies altamente correlacionados de um fator latente, um filtro de
correlação ingênuo pode escolher arbitrariamente um proxy e nunca convergir
no correto. A distinção "mesma base nunca filtra, base diferente filtra" é
a fronteira que evita esse problema por construção.

\section{\texttt{treinamento} --- Pedro\_Wise}

O "famigerado" Pedro\_Wise --- busca greedy multi-nível (forward/backward,
níveis 1 a 3) sobre GLM binomial (logística), otimizando KS (estatística de
Kolmogorov--Smirnov entre a distribuição de score de eventos e
não-eventos). Port fiel de um script R original, validado numericamente
contra o R (KS-dev bate em 10 casas decimais).

\subsection*{\texttt{treinar(df\_dev, df\_teste, ...)}}

Todas as colunas de \texttt{df\_dev}/\texttt{df\_teste} exceto \texttt{y}
são candidatas.

\begin{longtable}{P{4.3cm}P{1.8cm}P{1.4cm}P{6.2cm}}
\toprule
\textbf{Parâmetro} & \textbf{Tipo} & \textbf{Default} & \textbf{Explicação} \\
\midrule
\endhead
\texttt{criterio} & teste\,|\,dev\,|\,min & "teste" & \textbf{Não é cosmético} --- é o que forward/backward otimiza em toda a busca. O R original tem um bug que decide só por KS-dev; a validação do port mostrou que sob \texttt{criterio="dev"} o algoritmo aceita uma variável de puro ruído por acaso amostral. \texttt{"teste"} é o default recomendado para uso real. \\
\texttt{shadow\_probing} & bool & False & Reservado para um critério de parada alternativo (backlog) inspirado em Thomas, Hepp, Mayr \& Bischl (2017), \emph{Probing for Sparse and Fast Variable Selection with Model-Based Boosting}: aumentar as candidatas com cópias permutadas ("shadow variables") e parar quando uma shadow seria selecionada --- sinal de ajuste a ruído, sem precisar separar dev/teste para decidir o ponto de parada. \\
\texttt{forward\_simples} & bool & True & Nível 1: testa adicionar cada variável fora do modelo, uma por vez. \\
\texttt{transformacao\_simples\_nivel1} & bool & True & Nível 1: troca a versão de uma base já no modelo por outra versão da mesma base. \\
\texttt{backward\_simples\_nivel1} & bool & True & Nível 1: remove uma variável por vez, se o KS melhora. \\
\texttt{min\_vars\_para\_backward} & int & 5 & Backward só é tentado acima desse número de variáveis no modelo. \\
\texttt{forward\_duplo} & bool & True & Nível 2: testa pares de variáveis do top-\texttt{n\_best\_duplo} do forward simples --- pega sinergias que forward variável-a-variável não descobre sozinho. \\
\texttt{forward\_triplo} & bool & True & Nível 2.5: combina top-\texttt{n\_best\_triplo\_1} $\times$ \texttt{n\_best\_triplo\_2} do forward duplo em triplas. \\
\texttt{transformacao\_simples\_nivel2} & bool & True & Mesma troca de versão de base, repetida no nível 2. \\
\texttt{backward\_simples\_nivel2} & bool & True & Mesmo backward, repetido no nível 2. \\
\texttt{n\_best\_duplo} & int & 5 & Quantas variáveis do forward simples entram na busca de pares. \\
\texttt{n\_best\_triplo\_1}, \texttt{n\_best\_triplo\_2} & int & 3, 3 & Quantas variáveis de cada lado entram na busca de triplas. \\
\texttt{nivel3\_ativado} & bool & False & Ativa o backward complexo (nível 3): remove e re-roda o Pedro\_Wise inteiro a partir dali (recursão). Maior custo/risco combinatorial. \\
\texttt{n\_best\_backward} & int & 2 & Quantas variáveis avalia para remoção no nível 3. \\
\texttt{profundidade\_maxima\_nivel3} & int & 2 & Limite de profundidade da recursão do nível 3. \\
\texttt{p\_valor\_maximo} & float\,|\,None & None & Filtro opcional de significância --- não é o critério de seleção (o Pedro\_Wise nunca usa p-valor para decidir, sempre KS). \\
\bottomrule
\end{longtable}

Retorna \texttt{ResultadoTreinamento}: \texttt{variaveis}, \texttt{ks\_dev},
\texttt{ks\_teste}, \texttt{auc\_teste}, \texttt{taxa\_evento\_dev},
\texttt{taxa\_evento\_teste}, \texttt{coeficientes}, \texttt{estatisticas}
(diagnóstico pós-hoc, não influencia a seleção), \texttt{tabela\_decis}.

\section*{Bibliografia citada}
\begin{itemize}[leftmargin=1cm]
\item[] \label{friedman2008}\textbf{Friedman, J. H., \& Popescu, B. E.} (2008). \emph{Predictive learning via rule ensembles}. Annals of Applied Statistics. --- base de \texttt{esfera2}.
\item[] \label{kanter2015}\textbf{Kanter, J. M., \& Veeramachaneni, K.} (2015). \emph{Deep feature synthesis: Towards automating data science endeavors}. --- vocabulário de primitivas de \texttt{construcao}/\texttt{esfera1}.
\item[] \label{navaspalencia2020}\textbf{Navas-Palencia, G.} (2020). \emph{Optimal binning: mathematical programming formulation}. arXiv:2001.08025. --- inspiração de \texttt{bins\_monotonicos} em \texttt{categorizar}.
\item[] \textbf{Siddiqi, N.} \emph{Credit Risk Scorecards}. --- convenção de WOE em \texttt{categorizar}.
\item[] \label{meinshausen2010}\textbf{Meinshausen, N., \& Bühlmann, P.} (2008/2010). \emph{Stability Selection}. JRSS-B. --- motivação de \texttt{preselecao}.
\item[] \label{faletto2022}\textbf{Faletto, G., \& Bien, J.} (2022). \emph{Cluster Stability Selection}. --- desenho de \texttt{filtrar\_correlacao} (exceção para mesma base).
\item[] \label{thomas2017}\textbf{Thomas, J., Hepp, T., Mayr, A., \& Bischl, B.} (2017). \emph{Probing for Sparse and Fast Variable Selection with Model-Based Boosting}. Computational and Mathematical Methods in Medicine. --- \texttt{shadow\_probing} em \texttt{treinamento} (backlog).
\item[] \label{pargent2022}\textbf{Pargent, F., Pfisterer, F., Thomas, J., \& Bischl, B.} (2022). \emph{Regularized target encoding outperforms traditional methods...} --- nota de cautela em \texttt{transformar\_categoricas\_woe}.
\item[] \label{atkinson2021}\textbf{Atkinson, A. C., Riani, M., \& Corbellini, A.} (2021). \emph{The Box–Cox Transformation: Review and Extensions}. --- paralelo às transformações de potência de \texttt{categorizar}.
\end{itemize}

\vspace{1em}
\noindent Ver \texttt{docs/literatura/} no repositório para a lista completa
com links e sínteses mais longas por tópico.

\end{document}
